\section{Case Study and Evaluation}
\label{sec:zflean-case-study}

Internally, all constructions developed so far in the ZF model are \emph{sets}, including functions, relations, naturals, etc.
This section illustrates the intended \emph{user workflow} in \zflean{} on a classical result that exercises most of the infrastructure developed above, with an attempt to demonstrate how the interface is designed so that these implementation details rarely surface.
The chosen example is the \emph{currying isomorphism}, which is representative because it combines nested abstraction, function application and transformations, while remaining mathematically standard and proof-script friendly.

\subsection{Currying isomorphism}

\begin{theorem}[Currying isomorphism][isIso\_curry][ZFLean/Isomorphisms.html\#ZFSet.isIso_curry]
  \label{thm:isIso_curry}
  For any sets $A,B,C$, the following holds:
  \begin{equation*}
    C^{A\times B} \zbinop{\cong} (C^B)^A
  \end{equation*}
\end{theorem}
%
We use two mutually inverse functions to build the isomorphism, corresponding to the usual (un)currying operation. Those functions are defined as follows in \zflean.
%
\begin{definition}[Currying and uncurrying][ZFSet.currify, ZFSet.uncurrify][ZFLean/Isomorphisms.html\#ZFSet.currify, ZFLean/Isomorphisms.html\#ZFSet.uncurrify]
  We define set-level functions $\zunop{\mathsf{curry}} \colon \mathsf{ZFSet}$ and $\zunop{\mathsf{uncurry}} \colon \mathsf{ZFSet}$ as follows in \zflean, mirroring the mathematical definition and syntax:
  %
  \begin{equation*}
    \zunop{\mathsf{curry}} \triangleq \left(\mkern-5mu\zlambda{C^{A \times B}}{(C^B)^A}{f}{\mathcal{C}(f)}\mkern-5mu\right),\,
    \zunop{\mathsf{uncurry}} \triangleq \left(\mkern-5mu\zlambda{(C^B)^A}{C^{A \times B}}{f}{\mathcal{U}(f)}\mkern-5mu\right)
  \end{equation*}
  %
  where
  %
  \begin{equation*}
    \mathcal{C}(f) \triangleq \left(\zlambda{A}{C^B}{a}{\left(\zlambda{B}{C}{b}{\fapply{f}{\langle (a, b), \irrelprf \rangle}}\right)}\!\right)
  \end{equation*}
  %
  and
  %
  \begin{equation*}
    \mathcal{U}(f) \triangleq \left(\zlambda{A \times B}{C}{(a,b)}{\fapply{\left(\fapply{f}{\langle a, \irrelprf \rangle}\right)}{\langle b, \irrelprf \rangle}}\!\right)
  \end{equation*}
\end{definition}

We prove that $\zunop{\mathsf{curry}}$ and $\zunop{\mathsf{uncurry}}$ are total functions
on their respective domains, so that automation can pick those facts up when needed:
\[
  \mathsf{IsFunc}(C^{A \times B}, (C^B)^A, \zunop{\mathsf{curry}}) \quad\text{and}\quad \mathsf{IsFunc}((C^B)^A, C^{A \times B}, \zunop{\mathsf{uncurry}}).
\]

We then prove that these two functions are mutually inverse; this is the core of the proof of \cref{thm:isIso_curry}.
\begin{lemma}[][ZFSet.currify\_of\_uncurrify][ZFLean/Isomorphisms.html\#ZFSet.currify_of_uncurrify]
  \label{lem:curry-uncurry-inverse}
  \begin{equation*}
    \zunop{\mathsf{curry}} \zbinop{\circ} \zunop{\mathsf{uncurry}} = \Id{(C^B)^A} \quad \text{and} \quad
    \zunop{\mathsf{uncurry}} \zbinop{\circ} \zunop{\mathsf{curry}} = \Id{C^{A \times B}}
  \end{equation*}
\end{lemma}
\begin{proof}
  We prove only the first equality; the second follows a similar pattern.
  First, note that both sides contain side-conditions ensuring that everything is well-defined, all automatically discharged by the automation layer.
  We first reduce the equality of functions to a pointwise equality using function extensionality, using theorem \inlinelean|is_func_ext_iff| from \cref{ex:func-ext}.
  Then, let $f \in (C^B)^A$ be an arbitrary set-level curried function; we need to show the following equality:
  \begin{equation*}
    \fapply{(\zunop{\mathsf{curry}} \zbinop{\circ} \zunop{\mathsf{uncurry}})}{\langle f, \irrelprf \rangle} =
    \fapply{\Id{(C^B)^A}}{\langle f, \irrelprf \rangle}
  \end{equation*}
  The right-hand side reduces to $f$ by definition of the identity function and using the following theorem \leancodeptr{fapply\_Id}{ZFLean/Functions.html\#ZFSet.fapply_Id}:
  \begin{leanlisting}
theorem fapply_Id {A x : ZFSet} (hx : x ∈ A) : @ᶻ𝟙A ⟨x, $\irrelprf$⟩ = ⟨x, hx⟩
\end{leanlisting}
  On the left-hand side, we use that application distributes over composition---given by theorem {\small\leancodeptr{fapply\_composition}{ZFLean/Functions.html\#ZFSet.fapply_composition}}:
  \begin{leanlisting}
theorem fapply_composition {g f : ZFSet} {A B C : ZFSet}
  (hg : B.IsFunc C g) (hf : A.IsFunc B f) {x : ZFSet} (xA : x ∈ A) :
  @ᶻ(g ∘ᶻ f) ⟨x, $\irrelprf$⟩ = @ᶻg ⟨@ᶻf ⟨x, $\irrelprf$⟩, $\irrelprf$⟩
\end{leanlisting}
  This reduces the left-hand side to
  $\fapply{\zunop{\mathsf{curry}}}{\langle \fapply{\zunop{\mathsf{uncurry}}}{\langle f, \irrelprf \rangle}, \irrelprf \rangle}$.
  Unfolding the definitions of $\zunop{\mathsf{curry}}$ and $\zunop{\mathsf{uncurry}}$ and applying $\beta$-reduction twice from \cref{thm:lambda} then yields the following equality to prove:
  \begin{equation*}
    \mathcal{C}(\mathcal{U}(f)) = f
  \end{equation*}
  Again, we leverage the framework's functional extensionality theorem and apply it twice to reduce this equality to a pointwise one for any $a \in A$ and $b \in B$:
  \begin{equation*}
    \fapply{(\fapply{\mathcal{C}(\mathcal{U}(f))}{\langle a, \irrelprf \rangle})}{\langle b, \irrelprf \rangle} =
    \fapply{(\fapply{f}{\langle a, \irrelprf \rangle})}{\langle b, \irrelprf \rangle}
  \end{equation*}
  Further $\beta$-reductions are applied to the unfolded definition of $\mathcal{C}$, yielding:
  \begin{equation*}
    \fapply{\left(\zlambda{B}{C}{b'}{\fapply{\mathcal{U}(f)}{\langle (a, b'), \irrelprf \rangle}}\right)}{\langle b, \irrelprf \rangle} =
    \fapply{(\fapply{f}{\langle a, \irrelprf \rangle})}{\langle b, \irrelprf \rangle}
  \end{equation*}
  and to the unfolded definition of $\mathcal{U}$ as well, reducing the goal to the trivial equality
  \mbox{\(
    \fapply{\left(\fapply{f}{\langle a, \irrelprf \rangle}\right)}{\langle b, \irrelprf \rangle} =
    \fapply{(\fapply{f}{\langle a, \irrelprf \rangle})}{\langle b, \irrelprf \rangle}
    \)}.%
\end{proof}

The proof of \cref{thm:isIso_curry} then boils down to applying lemma \leancodeptr{isIso\_of\_two\_sided\_inverse}{ZFLean/Isomorphisms.html\#ZFSet.isIso_of_two_sided_inverse} with the equalities from \cref{lem:curry-uncurry-inverse}:
\begin{leanlisting}
theorem isIso_of_two_sided_inverse {A B : ZFSet} {f g : ZFSet}
  {hf : A.IsFunc B f} {hg : B.IsFunc A g}
  (left_inv : g ∘ᶻ f = 𝟙A) (right_inv : f ∘ᶻ g = 𝟙B) : A ≅ᶻ B
\end{leanlisting}

\begin{leanlisting}[float=t,
  title={Proof skeleton of the currying isomorphism.},
  label=lst:isIso-curry-skel]
theorem isIso_curry {A B C : ZFSet} : (A × B).funs C ≅ᶻ A.funs (B.funs C) := by
  have l_inv : uncurry ∘ᶻ curry = 𝟙((A × B).funs C) := $\irrelprf$
  have r_inv : curry ∘ᶻ uncurry = 𝟙(A.funs (B.funs C)) := by
    rw [is_func_ext_iff]
    intro f hf
    unfold curry uncurry
    rw [fapply_Id, fapply_composition, fapply_lambda $\irrelprf$ $\irrelprf$, fapply_lambda $\irrelprf$ $\irrelprf$]
    $\irrelprf$
  exact isIso_of_two_sided_inverse l_inv r_inv
\end{leanlisting}

All routine side-conditions (``is a partial/total function'', ``is a relation'') are discharged structurally and automatically by \inlinelean|zrel|/\inlinelean|zpfun|/\inlinelean|zfun| through automatic parameters and the appropriate attribute-tagged closure lemmas.
The skeleton of the proof script shown in \cref{lst:isIso-curry-skel} closely follows the mathematical argument described above.
At the proof-script level, the role of the automation layer is not to perform proof search, but to \emph{erase boilerplate}: the user writes the mathematical argument, and the automation layer discharges the ubiquitous well-formedness obligations generated by set-level definitions and rewriting.

\subsection{Size of the framework}
\label{subsec:artifacts}

\begin{table}[t]
  \centering
  \small
  \begin{tabular}{lrr}
    \hline
    \textbf{Component (modules)}                 & \textbf{LoC}   & \textbf{\#Theorems} \\
    \hline
    Relational calculus                          & 2,534          & 129                 \\

    Embeddings, isomorphisms                     & 1,626          & 27                  \\

    Other canonicals (Booleans, sums, rationals) & 1,362          & 128                 \\

    Naturals                                     & 1,259          & 180                 \\

    Integers                                     & 1,189          & 128                 \\

    Core glue and automation                     & 316            & 41                   \\
    \hline
    \textbf{Total}                               & \textbf{8,286} & \textbf{633}        \\
    \hline
  \end{tabular}
  \caption{Artifact sizes in lines of Lean code (LoC) and number of theorems, grouped by component.}
  \label{tab:artifacts-loc}
\end{table}

The artifact is a self-contained Lean~4 library layered on top of Mathlib's ZFC model consisting of roughly 8k lines of code (LoC), organized in a few modules summarized in \cref{tab:artifacts-loc}, and containing around 633 theorems.
The relational calculus is the largest component, representing about one third of the codebase.

